\section{Conclusion}
\label{sec:conclusion}

% NOTE: source arXiv 2505.22954 main.tex merges Conclusion and Limitations
% into a single "Conclusion and Limitations" section (sec:conclusion). This
% file and paper/sections/07_limitations.tex split that section's verbatim
% paragraphs by topic to match the template's separate Conclusion/Limitations
% slots. See paper/supplementary/source-attribution.md.

We introduce the Darwin G\"odel Machine (DGM), the first self-improving system powered by FMs with open-ended exploration, where progress on its evaluation benchmarks can directly translate into better self-improvement capabilities. We demonstrate the automatic discovery of better tools and FM systems, resulting in better performance on two benchmarks: SWE-bench and Polyglot. Through self-improvement and open-ended exploration, the DGM shows a continuous increase in performance, bringing us one step closer to self-accelerating, self-improving AI systems.

In conclusion, the DGM represents a significant step toward the automation of AI development through self-improving systems capable of editing their own codebase. While current limitations in compute and reasoning constrain its full potential, continued advances in FMs and infrastructure may unlock more powerful and general-purpose self-improvements. Provided that the safety concerns are carefully navigated (\Cref{sec:safety}), the future of self-improving AI systems and AI-Generating Algorithms~\citep{clune2019ai} holds immense promise to open-endedly evolve AI, continually rewriting or retraining itself in pursuit of greater capabilities aligned with human values.
